% Pursuing Stacks (cotes 134-2 to 134-8): archive cote 134-4, site page 49 (PDF page 50), typescript page 137.
% typescript. Draft compared: ps3.tex:2481-2539.
% AI-transcribed by Claude Opus 5.5 (claude-opus-5-5, adaptive thinking, effort high), Message Batches API, batch ps2 (msgbatch_01NVwr1Yuo8qJf5rykeGBXPK), 2026-11-12; not yet checked by eye.
\scanpage{134-4/49}
\mspage{137}

This expectation of $W$ to be recoverable in terms of the corresponding
homotopy structure $h = h_{W}$ on $M$ takes it's full meaning,  when joined with
another one,  \xout{which} namely that the latter can be canonically described in
terms of the category structure of $M$  and the corresponding notion of
0-connectedness.  This latter expectation is extremely strongly grounded,
and I'll come back to it circumstancially very soon I think  (I started
on it two days ago,  but then it got too late to take it to the end, and
yesterday was spent on some formal digressions ...)  The two ``expectations''
put together,  when realized by \xout{xxxxxxx}\inserted[typed]{carefully} cutting out the suitable notions,
\over{\unsure{would}}{should} imply that the structure of any ``canonical modelizer'' is indeed deter%
mined ``canonically'' in terms of it's category structure alone.

To come back to the relationship between test categories and catego-%
ries of the type $M^{\unsure{\circ|0}}_{c}$ ,  the idea which has been lurking \xout{xxxxx} lately is
that possibly,  test categories can be viewed as no more,  no less,  as
categories endowed with \xout{the (unique} a homotopy structure (necessarily
unique)  which is a contractibility structure,  and for which \uline{all ob\over{\unsure{g}}{j}ects
are contractible} .  At any rate,  there must be a very close relationship
between the two notions,  which I surely want to understand.  But as it
would be quite unreasonable to restrict the notion of \inserted{a} test categor\over{ies}{y}
(and it's weak and strong variants) to categories admitting finite pro-%
ducts,  this shows that for a satisfactory understanding of the above re-%
lationship,  we should be able to work with contractibility structures,
and presumably too with weak homotopy interval structures too,  in catego-%
ries $A$ where we do not make the assumption of stability under products.
Thus in the \xout{xxxxx}\inserted[typed]{outline} of the last few days,  I still was'nt general enough
it would seem !  The situation reminds me rather strongly of the early
stages when developing the language of sites,  and restricting to sites
where fiber products exist  -  this seemed by then a very weak and natural
assumption \deleted{indeed},  before it appeared (first to Giraud, I believe) that
it was quite an awkward and artificial restriction indeed,  which had to
be overcome in order to work really at ease ...

All this now gives a lot of interesting things to look up in the
short run!  I'll make a provisional plan of work, as follows:

a)  Relation between a homotopy structure and the $\hand{$\pi$}_{o}$ functor,  and
description of the so-called canonical homotopy structures.

b)  Write down in the end the ``key result'' on test functors $A$\hand{$\rightarrow$}(Cat),
which  is overripe since the reflections  of four days ago (par. 47 ) .

\margin{49}
